% ============================================================
%  2.2 – Nguyên lý thiết kế hình thức  (0.5 điểm)
%  Trạng thái: ĐÃ HOÀN THIỆN
% ============================================================

\subsection{Nguyên lý kiến trúc thông tin (Information Architecture)}

Trước khi đi vào thiết kế thị giác, nhóm áp dụng \textbf{5 nguyên lý IA} để tổ chức
và trình bày thông tin hiệu quả:

\begin{enumerate}
  \item \textbf{Objects (Đối tượng):} Mỗi thực thể trong hệ thống (job card, user
    profile, application) được thiết kế như một đối tượng độc lập với trạng thái và
    hành vi riêng. Job card hiển thị trạng thái ``Khẩn cấp'', ``Đã ứng tuyển'',
    ``Đã đóng'' bằng badge màu sắc phân biệt, không cần người dùng vào chi tiết.

  \item \textbf{Choices (Lựa chọn):} Giới hạn số hành động hiển thị trên mỗi màn hình.
    Trang JobDetail chỉ có 2 CTA chính: ``Ứng tuyển ngay'' (primary) và
    ``Bookmark'' (secondary); không bày nhiều nút cùng lúc để tránh cognitive overload.

  \item \textbf{Disclosure (Tiết lộ dần):} Thông tin được hiển thị theo từng lớp --
    job card chỉ tiết lộ tiêu đề, danh mục, mức lương và badge; khi nhấn vào mới
    hiện đầy đủ mô tả, kỹ năng yêu cầu và form ứng tuyển inline. Người dùng không
    bị ngợp thông tin ngay từ đầu.

  \item \textbf{Exemplars (Minh họa):} Mỗi danh mục công việc kèm icon trực quan
    (Lucide React: \texttt{Code2} cho IT, \texttt{Palette} cho Design,
    \texttt{BookOpen} cho Gia sư, \texttt{Camera} cho Media...) giúp nhận diện
    danh mục ngay mà không cần đọc tên text.

  \item \textbf{Front Doors (Đa điểm vào):} Cùng một job có thể tiếp cận từ nhiều
    nguồn: Hero section gợi ý, khối SuggestedJobs (AI matching), thanh tìm kiếm
    full-text, hoặc URL trực tiếp \texttt{/jobs/:id}. React Router DOM đảm bảo
    mọi đường dẫn đều hoạt động khi reload nhờ cấu hình rewrite về \texttt{index.html}.
\end{enumerate}

\subsection{Nguyên lý thị giác (CARP)}

Thiết kế hình thức của UniJob được xây dựng trên \textbf{4 nguyên lý cốt lõi} từ lý thuyết
thiết kế giao diện hiện đại:

\begin{enumerate}
  \item \textbf{Proximity (Gần kề):} Các thông tin liên quan được nhóm lại thành
    đơn vị trực quan (card). Mỗi job card chứa đầy đủ 6--8 thông tin cốt lõi trong một
    khối bao quanh bởi viền và shadow nhẹ, giúp mắt người dùng di chuyển tự nhiên.

  \item \textbf{Alignment (Căn chỉnh):} Tất cả nội dung căn trái trên desktop, căn giữa
    trên mobile. Lưới Tailwind CSS đảm bảo căn chỉnh pixel-perfect trên mọi viewport.

  \item \textbf{Contrast (Tương phản):} Màu chủ đạo xanh lá (\texttt{\#10b981}) trên nền
    trắng đạt tỷ lệ tương phản 3.1:1 (WCAG AA cho large text). Text body \texttt{\#111827}
    trên trắng đạt 16.1:1, vượt chuẩn WCAG AAA.

  \item \textbf{Repetition (Lặp lại):} Design tokens CSS (\texttt{--color-primary},
    \texttt{--color-border}...) được dùng xuyên suốt toàn bộ codebase, đảm bảo tính nhất
    quán khi thay đổi theme mà không cần sửa từng component.
\end{enumerate}

\subsection{Hệ thống Typography}

\begin{table}[H]
\centering
\caption{Cấu trúc typography trong UniJob}
\renewcommand{\arraystretch}{1.2}
\begin{tabularx}{\textwidth}{|p{3.5cm}|p{3.5cm}|X|}
\hline
\textbf{Mức} & \textbf{Class Tailwind} & \textbf{Sử dụng} \\
\hline
Display Heading & \texttt{text-4xl font-bold} & Hero headline trang chủ \\
\hline
H1 Section      & \texttt{text-2xl font-bold} & Tiêu đề trang (JobDetail, Profile) \\
\hline
H2 Card title   & \texttt{text-lg font-semibold} & Tiêu đề job card \\
\hline
Body            & \texttt{text-sm} (14px)     & Mô tả, nội dung form \\
\hline
Caption         & \texttt{text-xs text-gray-500} & Metadata: ngày, tag, badge \\
\hline
\end{tabularx}
\end{table}

Font stack mặc định: \texttt{"Segoe UI", Roboto, Arial, sans-serif} -- không nhúng custom
font để tối ưu hiệu suất tải trang.

\subsection{Hệ thống Component}

\begin{table}[H]
\centering
\caption{Hệ thống component UI của UniJob}
\renewcommand{\arraystretch}{1.2}
\begin{tabularx}{\textwidth}{|p{2.5cm}|p{3.5cm}|X|}
\hline
\rowcolor{emeraldrow}
\textbf{\color{emeraldhead}Component} &
\textbf{\color{emeraldhead}Biến thể / Loại} &
\textbf{\color{emeraldhead}Mô tả \& Ví dụ sử dụng} \\
\hline
\multirow{3}{*}{\textbf{Button}} &
  Primary (nền emerald) & CTA chính: ``Ứng tuyển ngay'', ``Tải xuống PDF'' \\
\cline{2-3}
& Outline (viền emerald) & Hành động thứ cấp: ``Chia sẻ LinkedIn'', ``Hủy'' \\
\cline{2-3}
& Ghost (không viền/nền) & Điều hướng nội bộ. Tất cả có \texttt{disabled} + loading spinner \\
\hline
\textbf{Card} &
  \texttt{rounded-2xl shadow-sm} &
  Đơn vị giao diện cơ bản. Hover: \texttt{shadow-md + border-emerald-200}.
  Padding \texttt{p-4}/\texttt{p-6} tùy context \\
\hline
\multirow{3}{*}{\textbf{Badge \& Tag}} &
  Category badge & \texttt{bg-emerald-100} -- phân loại danh mục công việc \\
\cline{2-3}
& Urgent badge & \texttt{bg-orange-100} + icon Zap -- task cần người gấp \\
\cline{2-3}
& Skill tag & \texttt{bg-gray-100} font nhỏ -- không cạnh tranh sự chú ý \\
\hline
\multirow{3}{*}{\textbf{Toast}} &
  Success (tick xanh lá) & Ứng tuyển thành công, lưu profile \\
\cline{2-3}
& Error (X đỏ) & Lỗi mạng, form validation fail \\
\cline{2-3}
& Info (icon ℹ) & Tính năng đang phát triển (LinkedIn share) \\
\hline
\end{tabularx}
\end{table}

\subsection{Micro-interactions và Animation}

\begin{table}[H]
\centering
\caption{Micro-interactions và animation trong UniJob}
\renewcommand{\arraystretch}{1.2}
\begin{tabularx}{\textwidth}{|p{3.5cm}|p{3.5cm}|X|}
\hline
\textbf{Hiệu ứng} & \textbf{Kỹ thuật} & \textbf{Mục đích} \\
\hline
Hero fade-in & \texttt{framer-motion}: \texttt{opacity 0→1, y 20→0} (0.6s) & Tạo ấn tượng khi vào trang chủ \\
\hline
Apply form slide & \texttt{transition-all} CSS khi toggle & Form ứng tuyển hiện/ẩn mượt, không giật \\
\hline
Loading spinner & Vòng quay \texttt{border-emerald} thay thế content & Thông báo đang fetch từ Firestore \\
\hline
Skeleton loading & \texttt{animate-pulse} trên block gray & Tránh Layout Shift (CLS) khi data chưa về \\
\hline
\end{tabularx}
\end{table}

\subsection{Layout và Spacing System}

\begin{table}[H]
\centering
\caption{Hệ thống layout và spacing của UniJob}
\renewcommand{\arraystretch}{1.2}
\begin{tabularx}{\textwidth}{|p{3.5cm}|p{4cm}|X|}
\hline
\textbf{Quy tắc} & \textbf{Giá trị} & \textbf{Áp dụng} \\
\hline
Spacing system & Bội số 4px: \texttt{gap-4} (16px), \texttt{gap-6} (24px), \texttt{p-8} (32px) & Toàn bộ layout \\
\hline
Max container width & \texttt{max-w-7xl} (1280px) & Trang danh sách việc làm \\
\hline
Max container width & \texttt{max-w-5xl} (1024px) & Trang chi tiết, form ứng tuyển \\
\hline
Grid -- Desktop & \texttt{grid-cols-3} & Job cards (3 cột) \\
\hline
Grid -- Mobile & \texttt{grid-cols-1} & Responsive, sidebar collapse \\
\hline
Grid -- CVExport & \texttt{grid-cols-5} (tỷ lệ 3:2) & Preview trái + Panel xuất file phải \\
\hline
\end{tabularx}
\end{table}
